\subsection{次数}

设 K 为一个交换域。对于 \(\mathrm{K}((\mathbf{X}_{i})_{i\in I})\) 的每个元素 r，存在 u、 \(v\in\mathrm{K}[(\mathbf{X}_{i})_{i\in I}]\) 使得 \(v\neq0\) 且 \(r=\frac{u}{v}\) 。关系式 \(\frac{u}{v}=\frac{u_{1}}{v_{1}}\) （其中 \(v\neq0\) , \(v_{1}\neq0\) ）等价于 \(uv_{1}=vu_{1}\) ；如果 \(r\neq0\) ，我们有 \(u\neq0\) 且 \(u_{1}\neq0\) ，因此 \(\deg u+\deg v_{1}=\deg v+\deg u_{1}\) （IV，第 9 页），或者也可以写成 \(\deg u-\deg v=\deg u_{1}-\deg v_{1}\) 。于是有理整数 \(\deg u-\deg v\) 仅依赖于 r；我们称之为 r 的次数，或 r 的总次数，并记作 \(\deg r\) 。我们约定写 \(\deg0=-\infty\) 。若 \(J\subset I\) ，我们同样可以定义关于 \(X_{j}\) （指标 \(j\in J\) ）的次数。当 r 是多项式时，这些概念与 IV 第 2 页所定义的一致。

命题 1. --- 设 r, s 为两个有理分式。

\begin{enumerate}
\def\labelenumi{(\roman{enumi})}
\tightlist
\item
  如果 \(\deg r \neq \deg s\) ，我们有
\end{enumerate}

\[
r + s \neq 0 \quad a n d \quad \deg (r + s) = \sup (\deg r, \deg s).
\]

如果 \(\deg r = \deg s\) ，我们有 \(\deg (r + s) \leqslant \deg r\) 。

\begin{enumerate}
\def\labelenumi{(\roman{enumi})}
\setcounter{enumi}{1}
\tightlist
\item
  我们有 \(\deg(rs) = \deg r + \deg s\) 。
\end{enumerate}

我们可以限于 \(r\) 与 \(s\) 都非零的情形。

我们写 \(r = \frac{u}{v}\) , \(s = \frac{w}{z}\) ，其中 \(u, v, w, z\) 是非零多项式。我们有 \(rs = \frac{uw}{vz}\) ，因此

\[
\deg (r s) = \deg (u w) - \deg (v z) = \deg u - \deg v + \deg w - \deg z = \quad = \deg r + \deg s.
\]

另一方面，我们有 \(r + s = \frac{uz + vw}{vz}\) 。假设 \(\deg r \neq \deg s\) ，换言之 \(\deg u + \deg z \neq \deg w + \deg v\) 。则 \(uz + wv \neq 0\) ，并且

\[
\begin{array}{r l} \deg (r + s) & = \deg (u z + w v) - \deg (v z) \\ & = \sup (\deg (u z), \deg (w v)) - \deg (v z) \\ & = \sup (\deg (u z) - \deg (v z), \deg (w v) - \deg (v z)) \\ & = \sup (\deg r, \deg s). \end{array}
\]

假设 \(\deg r = \deg s\) ，即 \(\deg u + \deg z = \deg w + \deg v\) 。若 \(r + s \neq 0\) ，则我们有

\[
\begin{array}{r l} \deg (r + s) & = \deg (u z + w v) - \deg (v z) \\ & \leqslant \deg (u z) - \deg (v z) = \deg r. \end{array}
\]

* 因此，映射 \(r \mapsto -\deg r\) 是域 \(K((X_i)_{i \in I})\) 上的一个离散赋值。*

